% SEGMENT OLP-0552-S01
\documentclass[../../../include/open-logic-section]{subfiles}

\begin{document}

\olfileid{sth}{ordinals}{vn}	
\olsection[ஃபான் நாய்மனின் கட்டமைப்பு]{ஃபான் நாய்மனின் வரிசையெண் கட்டமைப்பு}

\olref[sth][ordinals][iso]{thm:woalwayscomparable} ஒரு கருத்தைத்
தூண்டுகிறது. நன்கு வரிசைப்படுத்தல்களின் பதிலிகளாகச் செயல்பட, \emph{வரிசை
வகைகள்} எனப்படும் சில பொருள்களை நாம் அறிமுகப்படுத்தலாம். $\tuple{A, <}$
என்ற நன்கு வரிசைப்படுத்தலின் வரிசை வகையை $\ordtype{A, <}$ என்று எழுதினால்,
பின்வரும் இரண்டு கொள்கைகள் நிறைவடையும் என்று எதிர்பார்ப்போம்:
\begin{align*}
	\ordtype{A, <} = \ordtype{B, \lessdot} & 
	\text{ எனில், எனில் மட்டுமே, } \ordeq{\tuple{A, <}}{\tuple{B, \lessdot}}\\
	\ordtype{A, <} < \ordtype{B, \lessdot}&
	\text{ எனில், எனில் மட்டுமே, }\ordeq{\tuple{A, <}}{\tuple{B_b, \lessdot_b}}\text{ ஏதோ ஓர் }b \in B
\end{align*}
மேலும், இயல் எண்களைச் சில குறிப்பிட்ட கணங்களாக அறிமுகப்படுத்துவது போலவே,
வரிசை வகைகளையும் \emph{சில குறிப்பிட்ட கணங்களாக} அறிமுகப்படுத்தலாம் என்று
நம்பலாம்.

இதற்கான மிகப் பொதுவான வழியும் நாம் பின்பற்றவுள்ள அணுகுமுறையும், சில
\emph{நியம} நன்கு வரிசைப்படுத்தப்பட்ட கணங்கள் வழியாக இந்த வரிசை வகைகளை
வரையறுப்பதே. இந்த நியமக் கணங்களை முதலில் அறிமுகப்படுத்தியவர் ஃபான் நாய்மன்:

\begin{defn}
$(\forall x \in A)x \subseteq A$ எனில், அப்போது, அப்போது மட்டுமே,
$A$ என்ற கணம் \emph{கடப்புப் பண்புடையது}. மேலும் $A$ கடப்புப்
பண்புடையதாகவும் $\in$ ஆல் நன்கு வரிசைப்படுத்தப்பட்டதாகவும் இருந்தால்,
அப்போது, அப்போது மட்டுமே, $A$ ஒரு \emph{வரிசையெண்}.
\end{defn}
\noindent
இனி வரிசையெண்களுக்கு கிரேக்க எழுத்துகளைப் பயன்படுத்துவோம். $\alpha$ ஒரு
வரிசையெண் எனில், $\in_\alpha =
\Setabs{\tuple{x, y} \in \alpha^2}{x \in y}$ என அமையும்
$\tuple{\alpha, \in_\alpha}$ ஒரு நன்கு வரிசைப்படுத்தல் என்பது
வரையறையிலிருந்து உடனே பின்வருகிறது. ஆகவே, குறியீட்டைப் பயன்படுத்துவதில்
சிறிது தளர்வாக, $\alpha$ \emph{தானே} ஒரு நன்கு வரிசைப்படுத்தல் என்று
கூறலாம்.

நமது முதல் சில வரிசையெண்கள் இதோ:
\[
	\emptyset, \{\emptyset\}, 
	\{\emptyset, \{\emptyset\}\}, 
	\{\emptyset, \{\emptyset\}, \{\emptyset, \{\emptyset\}\}\}, \ldots
\]
நமது முடிவிலி அடிகோளில், அதாவது $\omega$ பற்றிய ஃபான் நாய்மனின்
வரையறையில், நாம் சந்தித்த முதல் சில வரிசையெண்களே இவையும் என்பதை நீங்கள்
கவனிப்பீர்கள் (\olref[sth][z][infinity-again]{sec} ஐப் பார்க்கவும்).
இது தற்செயல் அன்று. வரிசையெண்கள் பற்றிய ஃபான் நாய்மனின் வரையறை இயல்
எண்களை வரிசையெண்களாகக் கருதுகிறது; அதேநேரத்தில் முடிவிலிக்கடந்த
வரிசையெண்களையும் அனுமதிக்கிறது.

எப்போதும்போல இப்போதும் நாம் கேட்கலாம்: இவை \emph{உண்மையிலேயே}
வரிசையெண்களா? அல்லது வரிசையெண்களாக நாம் \emph{கருதக்கூடிய} சில கணங்களை
மட்டுமே ஃபான் நாய்மன் வழங்கியுள்ளாரா? இந்தக் கேள்வியைப் பற்றிய விவாதங்கள்,
\olref[sfr][rel][ref]{sec}, \olref[sfr][arith][ref]{sec},
\olref[sfr][infinite][dedekindsproof]{sec},
\olref[sth][z][nat]{sec} ஆகியவற்றில் நாம் நடத்திய விவாதங்களைப் போன்றவை;
எனவே இந்தக் கருத்தை மேலும் வலியுறுத்தமாட்டோம். இனி ``வரிசையெண்கள்'' என்று
கூறும்போது ``ஃபான் நாய்மன் வரிசையெண்களையே'' குறிப்போம்.

\end{document}
